\chapter{MEMORIAS DE CÁLCULO Y TRAZABILIDAD}
\label{anx:calculos}

\section{Corrida estructural R01}

La fuente primaria es el libro \texttt{evidencias/N3\_R01/R01\_SAP\_ORIGINAL.xlsx}. La Tabla~\ref{tab:traza-r01} identifica las tablas empleadas.

\begin{longtable}{P{3.5cm}P{4.0cm}P{5.0cm}}
\caption{Procedencia tabular de las variables R01}\label{tab:traza-r01}\\
\toprule
\textbf{Variable} & \textbf{Hoja} & \textbf{Uso en el informe} \\
\midrule
Unidades, versión y norma & Program Control & Identidad N--mm--C, SAP2000 24.2.0 y AISC 360-16 \\
Sección asignada & Frame Section Assignments & Confirmación interna W6X9 en 20 elementos de barra \\
Nodos y barras & Joint Coordinates; Connectivity - Frame & Geometría y conectividad R01 \\
Cargas & Joint Loads - Force & Tres fuerzas F3 de $-9948$ N \\
Casos y combinaciones & Load Case Definitions; Combination Definitions & Reconstrucción de SLS y ULS \\
Peso propio & Load Pattern Definitions & Multiplicadores 1,0 y 0,0 \\
Desplazamientos & Joint Displacements & U1, U2, U3 y flecha relativa \\
Fuerzas internas & Element Forces - Frames & P, V2, V3, T, M2 y M3 \\
Reacciones & Base Reactions; Joint Reactions & Equilibrio F1--F3 y M1--M3 \\
Diseño & Steel Sum - AISC 360-16 & D/C y combinación gobernante \\
\bottomrule
\end{longtable}

La exportación permite comparar los resultados del caso analizado. Las definiciones completas de secciones y materiales, junto con las opciones del análisis, se conservan en el archivo nativo \nolinkurl{evidencias/N3_R01/V5_2_N3_W6X9_R01.sdb}. Las tablas exportadas y el archivo nativo constituyen el respaldo del modelo estructural.

\subsection{Material y propiedades mecánicas del bastidor}
\label{sec:material-trazabilidad}

Para el bastidor de perfiles W6×9 se consideró acero ASTM A36 y se realizó la simulación estructural en SAP2000 con las propiedades mecánicas de este material. Esta especificación corresponde al bastidor; los centradores, ejes, rieles y demás componentes se dimensionaron con los materiales indicados en sus respectivos apartados.

La tabla \emph{Frame Section Assignments} registra W6X9 como sección de análisis y de diseño en los 20 elementos de barra. El campo \texttt{MatProp = Default} indica que se utiliza el material asociado a la sección, salvo una sobrescritura específica \cite{csi_material_overwrites}.

Las propiedades mecánicas empleadas fueron $E=199947{,}979$ MPa, $F_y=248{,}211$ MPa y $F_u=399{,}896$ MPa, según la comprobación de la barra 15 presentada en \nolinkurl{evidencias/N3_R01/04_comprobacion_uls_barra15.png}. Los resultados corresponden a las cargas, los apoyos y las idealizaciones del modelo de barras; su alcance no incluye la verificación de uniones, contactos ni del funcionamiento de la máquina completa.

\subsection{Respaldo del cálculo manual y sus extensiones}
\label{sec:respaldo-complementario}

El desarrollo manual del capítulo VI se contrasta con las hojas de coordenadas, cargas, reacciones y desplazamientos de \nolinkurl{evidencias/N3_R01/R01_SAP_ORIGINAL.xlsx}. No se dispone del libro complementario \nolinkurl{calculos/CALCULOS_SOPORTE_N3_N4_N5.xlsx}, citado como soporte de la hoja \emph{Contacto N3}; por ello, la extensión automática de ese desarrollo no forma parte de los entregables verificables. El libro mecánico R05 contiene las hojas \emph{Ruedas}, \emph{PAD\_Contacto} y \emph{Auditoria\_R01}, pero responde a sus propios casos y no sustituye esa hoja ni recalcula las corridas de cuatro y cinco PAD.

\section{Resultados de control R01}

\begin{table}[H]
\centering
\small
\caption{Resultados de control preservados}
\begin{tabularx}{\textwidth}{YP{4.0cm}}
\toprule
\textbf{Dato} & \textbf{Valor} \\
\midrule
Geometría de ejes & 1860×1900 mm \\
Carga externa total & 29,844 kN \\
Peso propio exportado & 1,64874 kN \\
SLS máximo nodal $|U3|$ & 1,353133 mm \\
SLS flecha relativa nodo 7 & 1,235802 mm respecto de cuerda 3--14 \\
ULS D/C máximo & 0,277132, barras 12, 13, 15 y 16 \\
\bottomrule
\end{tabularx}
\end{table}

\section{Cálculo mecánico R05 y aspersión I01}

La memoria mecánica R05 contiene los cálculos de integración de la mesa, mientras que la memoria de aspersión I01 reúne las ecuaciones y los parámetros hidráulicos. Estos desarrollos complementan el análisis estructural R01 y se presentan por separado para distinguir las hipótesis y los resultados de cada subsistema.

\subsection{Procedencia de las cargas de ruedas y riel}
\label{sec:traza-ruedas-riel}

Las cargas de ruedas y riel del apartado~\ref{sec:ruedas-rieles} se obtuvieron mediante el cálculo mecánico R05. La Tabla~\ref{tab:traza-ruedas-riel} relaciona cada magnitud con su hoja y celda en el libro \nolinkurl{evidencias/N3_R05/02_CALCULOS_EDITABLES_N3_R05.xlsx}. El procedimiento se desarrolla en \nolinkurl{evidencias/N3_R05/01_CALCULO_R05_REPRODUCIBLE.py} y los resultados se almacenan en \nolinkurl{evidencias/N3_R05/19_RESULTADOS_NUMERICOS_R05.json}. Este cálculo es distinto de la simulación estructural R01.

\begingroup
\small\singlespacing
\begin{longtable}{P{3.4cm}P{4.0cm}P{5.1cm}}
\caption{Trazabilidad de ruedas y riel en el libro mecánico R05}\label{tab:traza-ruedas-riel}\\
\toprule
\textbf{Magnitud} & \textbf{Hoja y celda} & \textbf{Valor y alcance} \\
\midrule
Masa de cálculo & \texttt{Entradas!C7} & 3500 kg; escenario de diseño \\
Masa estimada & \texttt{Masa\_CG!B20} & 3439,831288 kg; balance de masas \\
Peso de cálculo & \texttt{Ruedas!B19} & 34,335 kN \\
Carga de servicio & \texttt{Ruedas!B7} & 21,865488 kN; factor 1,15 \\
Carga resistente local & \texttt{Ruedas!B8} & 26,618855 kN; factor 1,40 \\
Reacción de instalación & \makecell[l]{\texttt{Riel\_Envolvente}\\\texttt{B13}} & 26,741372 kN; apoyo del riel \\
Momento del vano simple & \makecell[l]{\texttt{Riel\_Envolvente}\\\texttt{B16}} & 3956529,787 N\,mm; carga resistente \\
Esfuerzos de cabeza y base & \makecell[l]{\texttt{Riel\_Envolvente}\\\texttt{B17:B18}} & 146,538140 y 95,338067 MPa \\
Flecha de servicio & \makecell[l]{\texttt{Riel\_Envolvente}\\\texttt{B19}} & 0,531864 mm; carga de servicio \\
\bottomrule
\end{longtable}
\endgroup

El caso manual se identifica como \texttt{3C\_W1\_MP} en el código y en \texttt{wheel\_cases} del JSON. El libro identifica primero el caso simétrico \texttt{3C-03}, con rueda 3 ausente y $e_x=e_y=+100$ mm; ambos alcanzan el mismo máximo sin factor de 19,013468 kN. Las funciones \texttt{wheel\_cases} y \texttt{rail\_envelope} contienen los barridos restantes. El reparto de cuatro contactos requiere comprobar su compatibilidad con las deformaciones y rigideces del conjunto.

El bloque \texttt{rail.governing\_continuous} del JSON localiza la reacción de instalación en el riel derecho, apoyo 2, caso \texttt{3C\_W4\_PM}, con centro de mesa en $Y=-3320{,}038$ mm y una carga de rueda de 26,618855 kN situada en $Y=-4042{,}038$ mm. El mínimo continuo, $-3{,}691891$ kN, corresponde a los apoyos de instalación riel--base, no al contacto rueda--riel. La sensibilidad continua usa restricciones bilaterales y no valida la retención ni el diseño civil.

\section{Correspondencia de nombres de cálculo y archivos}
La Tabla~\ref{tab:correspondencia-sap} relaciona los identificadores que aparecen en los archivos de respaldo con la denominación que se emplea en el informe.

\begingroup
\small\singlespacing
\setlength{\parskip}{0pt}
\setlength{\tabcolsep}{4pt}
\renewcommand{\arraystretch}{1.1}
\begin{longtable}{@{}>{\raggedright\arraybackslash}p{3.6cm}>{\raggedright\arraybackslash}p{\dimexpr(\textwidth-3.6cm-2\tabcolsep)/1\relax}@{}}
\caption{Correspondencia entre identificadores originales y denominaciones del informe}\label{tab:correspondencia-sap}\\
\toprule
\textbf{Denominación} & \textbf{Significado}\\

\midrule\endfirsthead
\caption[]{Correspondencia de identificadores y denominaciones (continuación)}\\
\toprule
\textbf{Denominación} & \textbf{Significado}\\

\midrule\endhead
\midrule\multicolumn{2}{r}{\textit{Continúa en la página siguiente.}}\\\endfoot
\bottomrule\endlastfoot
N3, N4, N5 & Bastidor con tres, cuatro o cinco apoyos sectoriales, respectivamente. Los apoyos reciben el rodamiento y son distintos de las cuatro ruedas del conjunto móvil.\\
R01 & Análisis estructural original realizado en SAP2000.\\
R03, R04, R05 & Revisiones mecánicas: reserva de masa y desarrollo de la integración. R05 identifica la revisión mecánica posterior al análisis original, no cinco apoyos.\\
R06, R06.1 & Integración de la mesa con el aplicador, los rieles y el drenaje (R06) y desarrollo de los centradores K2-M (R06.1). Configuración candidata vigente.\\
\texttt{LC\_G\_FRAME} & Peso propio del bastidor.\\
\texttt{LC\_G\_EXTERNAL} & Carga vertical externa aplicada en los apoyos sectoriales.\\
\texttt{SLS\_NOMINAL} & Combinación nominal de servicio: peso propio y carga externa con factor 1,00.\\
\texttt{SLS\_UNC\_125} & Servicio con carga externa amplificada por 1,25; peso propio con factor 1,00. Este factor no sustituye el factor mecánico 1,15.\\
\texttt{ULS\_1\_4D} & Combinación resistente: factor 1,40 sobre el peso propio y la carga externa, conforme a la Tabla~\ref{tab:combinaciones-r01}.\\
\texttt{DSTL1} & Identificador de la combinación resistente reportada para cuatro y cinco apoyos. Su equivalencia con la combinación del bastidor con tres apoyos no está demostrada.\\
\texttt{Frame}, \texttt{Frames} & Barra o elementos de barra del modelo.\\
\texttt{MatProp = Default} & Campo de asignación del material descrito en el apartado~\ref{sec:material-trazabilidad}. Para el bastidor W6×9 se consideró acero ASTM A36.\\
D/C & Razón demanda/capacidad; no es un factor de seguridad global.\\
PMM & Interacción de carga axial y flexión biaxial.\\
V5.2, I01 & Identificadores de la arquitectura y del cálculo de aspersión documentados.\\
PAD y liner & Apoyo sectorial del rodamiento y revestimiento de su superficie de contacto. En el informe se emplea «apoyos sectoriales».\\
K2-M, v5.7.3, v5.7.4 & Centrador de brazo pivotante de eje vertical derivado del concepto K2. La versión 5.7.3 es la vigente; la 5.7.4 es un estudio de refuerzo no adoptado.\\
U1, U2, U3 & Unidades de centrado: U1 y U3 son las referencias y U2 el empujador.\\
LOAD, WASH & Receptores del pasador de bloqueo en las estaciones de carga y de lavado.\\
E1, E2 & Etapas de limpieza: E1, con el rodamiento montado sobre la mesa, y E2, con sus elementos desmontados.\\
Libro I01 & Identificador documental de la hoja de cálculo hidráulico. En el texto se utiliza «hoja de cálculo hidráulico».\\
\end{longtable}

\endgroup
